% Auto-generated complete chapter from 21_twinsight_demo_video_script_and_shotlist.md
\chapter{21\_twinsight\_demo\_video\_script\_and\_shotlist}
\label{complete:root_024}

\begin{sourcemeta}
\textbf{Fuente:} \href{file:///E:/Laboral/21_twinsight_demo_video_script_and_shotlist.md}{21\_twinsight\_demo\_video\_script\_and\_shotlist.md}\\
\textbf{Seccion:} root\\
\textbf{Estado:} canonical\\
\textbf{Rol:} Modulo principal\\
\textbf{Lineas originales:} 912\\
\textbf{Ultima modificacion:} 2026-06-11 15:18\\
\textbf{Deduplicacion conservadora:} 0 bloques repetidos omitidos; 0 caracteres repetidos no reimpresos.
\end{sourcemeta}

\section*{Resumen de orientacion}
This document defines the video plan for presenting \textbf{TwinSight X500} as the flagship portfolio asset.

\section*{Contenido completo depurado}
\addcontentsline{toc}{section}{Contenido completo depurado}




\subsection{Purpose}




This document defines the video plan for presenting \textbf{TwinSight X500} as the flagship portfolio asset.




It includes:




\begin{itemize}
\item 60-second demo script;
\item 90-second demo script;
\item 120-second demo script;
\item silent version;
\item LinkedIn cutdown;
\item technical breakdown version;
\item shot list;
\item capture checklist;
\item captions;
\item voiceover text;
\item editing structure;
\item export formats;
\item quality gate.
\end{itemize}




The goal is to make TwinSight understandable quickly to recruiters while still giving technical reviewers enough evidence to care.




\medskip\hrule\medskip




\section{1. Video objective}




The demo video must communicate, in under two minutes:




\begin{enumerate}
\item what TwinSight X500 is;
\item why it exists;
\item what problem it solves;
\item what features were built;
\item what technical skills it demonstrates;
\item what the measurable outcomes were;
\item what role Alexander had in the project.
\end{enumerate}




The video should not look like a thesis lecture. It should look like a concise technical product demo.




\medskip\hrule\medskip




\section{2. Primary audience}




\subsection{Recruiter audience}




Recruiters need to understand:




\begin{itemize}
\item this is a real Unity/WebGL project;
\item the candidate can build interactive systems;
\item the project is relevant to technical visualization, simulation, XR, digital twin, Unity and WebGL roles;
\item there is enough substance to justify an interview.
\end{itemize}




\subsection{Technical reviewer audience}




Technical reviewers need to see:




\begin{itemize}
\item Unity runtime interaction;
\item asset optimization;
\item inspection systems;
\item cross-section/clipping;
\item visual modes;
\item UI behavior;
\item WebGL/browser deployment;
\item evaluation metrics;
\item limitations handled honestly.
\end{itemize}




\medskip\hrule\medskip




\section{3. Core message}




Use this as the video thesis:




\begin{mdcode}
Bloque de codigo: text
TwinSight X500 transforms drone assembly documentation and CAD-derived assets into an optimized Unity WebGL technical visualization prototype for interactive inspection and spatial understanding.
\end{mdcode}




Shorter version:




\begin{mdcode}
Bloque de codigo: text
Seeing parts is not enough. Users need to understand spatial relationships.
\end{mdcode}




\medskip\hrule\medskip




\section{4. Video versions to produce}




\subsection{Required versions}




\begin{center}
\scriptsize
\begin{longtable}{>{\raggedright\arraybackslash}p{0.305\linewidth} >{\raggedright\arraybackslash}p{0.305\linewidth} >{\raggedright\arraybackslash}p{0.305\linewidth}}
\toprule
Version & Duration & Use \\
\midrule
\endfirsthead
\toprule
Version & Duration & Use \\
\midrule
\endhead
Recruiter demo & 60 s & LinkedIn, quick applications \\
Portfolio demo & 90 s & Portfolio hero/project page \\
Technical breakdown & 120 s & GitHub, technical reviewers \\
Silent captioned version & 60–90 s & LinkedIn autoplay \\
Vertical cutdown & 30–45 s & Optional social post \\
\bottomrule
\end{longtable}
\end{center}




\subsection{Recommended production order}




\begin{enumerate}
\item 90-second portfolio demo.
\item 60-second recruiter cut.
\item 120-second technical breakdown.
\item Silent LinkedIn version.
\end{enumerate}




Do not start with the 120-second version. It can become too dense.




\medskip\hrule\medskip




\section{5. Visual style}




\subsection{Recommended style}




\begin{center}
\scriptsize
\begin{longtable}{>{\raggedright\arraybackslash}p{0.455\linewidth} >{\raggedright\arraybackslash}p{0.455\linewidth}}
\toprule
Element & Direction \\
\midrule
\endfirsthead
\toprule
Element & Direction \\
\midrule
\endhead
Background & Dark or neutral \\
Text & Minimal, large, readable \\
Pace & Fast but not chaotic \\
Transitions & Simple cuts / fades \\
Captions & Short phrases \\
Audio & Clean voiceover or no voice \\
Music & Optional, subtle, non-distracting \\
Visual priority & App footage first \\
Branding & Minimal \\
\bottomrule
\end{longtable}
\end{center}




\subsection{Avoid}




\begin{itemize}
\item long intro animation;
\item thesis defense tone;
\item generic motivational music;
\item excessive text;
\item tiny UI text;
\item shaky screen recording;
\item uncompressed laggy capture;
\item showing development errors;
\item overexplaining AI;
\item putting legal disclaimers in the main video.
\end{itemize}




\medskip\hrule\medskip




\section{6. Capture requirements}




\subsection{Minimum capture quality}




\begin{center}
\scriptsize
\begin{longtable}{>{\raggedright\arraybackslash}p{0.455\linewidth} >{\raggedright\arraybackslash}p{0.455\linewidth}}
\toprule
Parameter & Recommendation \\
\midrule
\endfirsthead
\toprule
Parameter & Recommendation \\
\midrule
\endhead
Resolution & 1920×1080 \\
Frame rate & 30 fps minimum \\
Format & MP4 H.264 \\
Audio & 48 kHz if voiceover \\
Cursor & Visible only if useful \\
UI scale & Large enough \\
Browser & Clean window \\
Notifications & Off \\
Desktop clutter & Hidden \\
\bottomrule
\end{longtable}
\end{center}




\subsection{Capture tools}




Use any stable screen recorder:




\begin{itemize}
\item OBS Studio;
\item Nvidia ShadowPlay;
\item Windows Game Bar;
\item macOS screen recording;
\item Unity Recorder if recording inside editor;
\item browser capture if demo is live.
\end{itemize}




\subsection{Capture rule}




Record more footage than needed. Edit down later.




\medskip\hrule\medskip




\section{7. Shot list overview}




\subsection{Required shots}




\begin{center}
\scriptsize
\begin{longtable}{>{\raggedright\arraybackslash}p{0.455\linewidth} >{\raggedright\arraybackslash}p{0.455\linewidth}}
\toprule
Shot & Purpose \\
\midrule
\endfirsthead
\toprule
Shot & Purpose \\
\midrule
\endhead
Hero full drone & Immediate project recognition \\
Problem/2D reference & Explain why project exists \\
Orbit/inspect & Show real-time 3D \\
Component selection & Show interaction \\
Info panel & Show technical UI \\
Exploded view & Show spatial relationships \\
Cross-section/clipping & Show technical visualization \\
Visual modes & Show rendering/inspection systems \\
Optimization before/after & Show technical asset work \\
Metrics/evaluation & Show validation \\
Closing links & Call to action \\
\bottomrule
\end{longtable}
\end{center}




\subsection{Optional shots}




\begin{center}
\scriptsize
\begin{longtable}{>{\raggedright\arraybackslash}p{0.455\linewidth} >{\raggedright\arraybackslash}p{0.455\linewidth}}
\toprule
Shot & Purpose \\
\midrule
\endfirsthead
\toprule
Shot & Purpose \\
\midrule
\endhead
Mobile layout & Web/mobile proof \\
Wireframe overlay & Technical art proof \\
CAD mesh vs optimized mesh & Pipeline proof \\
Unity editor hierarchy & Implementation proof \\
Code snippet & Technical credibility \\
ARA not included & Avoid dilution \\
\bottomrule
\end{longtable}
\end{center}




\medskip\hrule\medskip




\section{8. 60-second recruiter demo script}




\subsection{Structure}




\begin{center}
\scriptsize
\begin{longtable}{>{\raggedright\arraybackslash}p{0.305\linewidth} >{\raggedright\arraybackslash}p{0.305\linewidth} >{\raggedright\arraybackslash}p{0.305\linewidth}}
\toprule
Time & Visual & Caption / Voice \\
\midrule
\endfirsthead
\toprule
Time & Visual & Caption / Voice \\
\midrule
\endhead
0–5 s & Full drone hero & TwinSight X500 — Unity WebGL technical visualization \\
5–10 s & 2D/manual vs 3D viewer & Drone assembly documentation can be hard to understand spatially \\
10–18 s & Orbit around drone & I built an interactive 3D inspection prototype in Unity \\
18–26 s & Component selection & Users can select components and inspect technical information \\
26–34 s & Exploded view & Exploded view reveals spatial relationships between parts \\
34–42 s & Cross-section & Cross-section tools support internal inspection \\
42–50 s & Visual modes & Visual modes help analyze structure, visibility and form \\
50–56 s & Optimization metric & CAD-derived geometry optimized for WebGL deployment \\
56–60 s & Closing screen & Unity · C\# · WebGL · Blender · Technical Visualization \\
\bottomrule
\end{longtable}
\end{center}




\subsection{60-second voiceover}




\begin{mdcode}
Bloque de codigo: text
This is TwinSight X500, a Unity WebGL technical visualization prototype for inspecting the assembly of a Holybro X500 V2 drone.

The project addresses a common problem in technical documentation: 2D diagrams can show parts, but users still have to reconstruct spatial relationships mentally.

I built an interactive 3D inspection experience with component selection, technical information panels, exploded view, cross-section tools and multiple visual modes.

The asset pipeline converts CAD-derived geometry into optimized real-time assets for browser deployment.

TwinSight was developed as my Multimedia Engineering thesis and evaluated using usability and workload methods including SUS, NASA-TLX Raw and Think-Aloud.

My role covered Unity/C# implementation, runtime interaction systems, technical UI, Blender-based asset preparation, documentation and final integration.
\end{mdcode}




\subsection{60-second captions}




Use these captions:




\begin{mdcode}
Bloque de codigo: text
TwinSight X500
Unity WebGL technical visualization

Problem:
2D documentation requires mental reconstruction

Interactive 3D inspection
Unity · C# · WebGL

Component selection
Technical information panels

Exploded view
Spatial relationships

Cross-section tools
Internal inspection

Visual modes
Real-time analysis

CAD-to-realtime optimization
Browser deployment

Built as Multimedia Engineering thesis
Available for remote Unity / real-time 3D work
\end{mdcode}




\medskip\hrule\medskip




\section{9. 90-second portfolio demo script}




\subsection{Structure}




\begin{center}
\scriptsize
\begin{longtable}{>{\raggedright\arraybackslash}p{0.305\linewidth} >{\raggedright\arraybackslash}p{0.305\linewidth} >{\raggedright\arraybackslash}p{0.305\linewidth}}
\toprule
Time & Visual & Message \\
\midrule
\endfirsthead
\toprule
Time & Visual & Message \\
\midrule
\endhead
0–5 s & Title + drone hero & What the project is \\
5–15 s & 2D problem / static diagram & Why it exists \\
15–25 s & Full interactive viewer & Unity WebGL prototype \\
25–35 s & Component selection & Interaction systems \\
35–45 s & Technical UI & Information architecture \\
45–55 s & Exploded view & Assembly relationships \\
55–65 s & Cross-section/clipping & Inspection tools \\
65–75 s & Visual modes & Rendering/analysis modes \\
75–82 s & Optimization stats & CAD-to-realtime proof \\
82–88 s & Evaluation metrics & Usability/workload proof \\
88–90 s & Links/CTA & Portfolio/GitHub/contact \\
\bottomrule
\end{longtable}
\end{center}




\subsection{90-second voiceover}




\begin{mdcode}
Bloque de codigo: text
TwinSight X500 is a Unity WebGL technical visualization prototype for inspecting and understanding the assembly of a Holybro X500 V2 drone.

The starting problem was technical documentation. Static diagrams and manuals can identify parts, but they often force users to mentally reconstruct how components relate in space.

TwinSight turns that assembly context into an interactive real-time 3D viewer.

Users can orbit the drone, select components, highlight parts and read technical information through the UI.

The exploded view separates components while preserving their spatial relationship, making the structure easier to inspect.

Cross-section and clipping tools allow users to analyze internal areas that would be difficult to understand from external views alone.

The prototype also includes multiple visual modes, including realistic, ghosted, blueprint, solid color and wireframe-style inspection modes.

A major part of the work was the CAD-to-realtime pipeline. CAD-derived geometry had to be cleaned, optimized and prepared in Blender before integration into Unity WebGL.

The project was evaluated using usability and workload methods such as SUS, NASA-TLX Raw and Think-Aloud.

My work covered Unity and C# implementation, runtime interaction systems, technical UI, Blender-based asset preparation, documentation, evaluation design and final integration.

The result is an academic prototype that demonstrates real-time 3D, WebGL deployment, technical visualization and CAD-to-realtime optimization.
\end{mdcode}




\subsection{90-second captions}




\begin{mdcode}
Bloque de codigo: text
TwinSight X500
Unity WebGL drone assembly visualization

Problem:
2D diagrams show parts, but not always spatial relationships

Solution:
Interactive real-time 3D inspection

Feature:
Component selection + technical UI

Feature:
Exploded view for assembly structure

Feature:
Cross-section / clipping tools

Feature:
Visual inspection modes

Pipeline:
CAD-derived assets → optimized real-time model

Evaluation:
SUS · NASA-TLX Raw · Think-Aloud

Role:
Unity/C# · WebGL · Blender · Technical Visualization

Portfolio / GitHub / Contact
\end{mdcode}




\medskip\hrule\medskip




\section{10. 120-second technical breakdown script}




\subsection{Structure}




\begin{center}
\scriptsize
\begin{longtable}{>{\raggedright\arraybackslash}p{0.305\linewidth} >{\raggedright\arraybackslash}p{0.305\linewidth} >{\raggedright\arraybackslash}p{0.305\linewidth}}
\toprule
Time & Visual & Message \\
\midrule
\endfirsthead
\toprule
Time & Visual & Message \\
\midrule
\endhead
0–6 s & Title and hero & Project identity \\
6–18 s & Problem / documentation & Research motivation \\
18–30 s & App overview & Unity WebGL viewer \\
30–42 s & Selection system & Runtime interaction \\
42–54 s & UI panels & Technical data display \\
54–66 s & Exploded view & Assembly relation system \\
66–78 s & Cross-section & Inspection/clipping \\
78–90 s & Visual modes & Rendering modes \\
90–102 s & CAD pipeline & Asset optimization \\
102–112 s & Metrics & Reduction/evaluation \\
112–120 s & Role + closing & Candidate contribution \\
\bottomrule
\end{longtable}
\end{center}




\subsection{120-second voiceover}




\begin{mdcode}
Bloque de codigo: text
This is TwinSight X500, a Unity WebGL technical visualization prototype built to inspect and understand the assembly of a Holybro X500 V2 drone.

The research problem comes from technical assembly documentation. Static diagrams and manuals identify components, but they often require users to mentally reconstruct spatial relationships between parts, subassemblies and fasteners.

TwinSight explores an interactive alternative: a browser-accessible real-time 3D viewer focused on inspection and spatial understanding.

The prototype allows users to orbit the drone, select components, highlight parts and access technical information panels.

The component selection system is designed around inspection workflows. The goal is not just to click a mesh, but to connect geometry, identity and technical information.

The exploded view separates the drone into readable groups while preserving assembly logic. This helps users understand how components relate to each other in space.

Cross-section and clipping tools support internal inspection, allowing hidden areas and layered components to be analyzed more clearly.

The project also includes visual modes such as realistic, ghosted, blueprint, solid color and wireframe-style views. These modes support different inspection tasks.

A major technical part of the project was the CAD-to-realtime pipeline. Dense CAD-derived geometry had to be converted, cleaned, optimized and prepared for Unity WebGL deployment.

The workflow included Blender-based cleanup, retopology and asset preparation, with the goal of preserving component identity while reducing the model enough for browser-based interaction.

Reported thesis metrics include geometry reduction from over 6.5 million triangles to approximately 95,617 optimized triangles, plus usability and workload evaluation using SUS, NASA-TLX Raw and Think-Aloud. These final values should be verified against the thesis report before public release.

My role covered Unity/C# implementation, runtime interaction systems, technical UI, CAD-to-realtime asset preparation, documentation, evaluation design and final integration.

TwinSight is an academic prototype, not a commercial engineering tool, but it demonstrates the kind of real-time 3D technical visualization work I want to continue building.
\end{mdcode}




\subsection{120-second captions}




\begin{mdcode}
Bloque de codigo: text
TwinSight X500
Unity WebGL technical visualization prototype

Problem:
Assembly documentation can hide spatial relationships

Goal:
Interactive 3D inspection and understanding

Runtime system:
Component selection and highlighting

Technical UI:
Part information panels

Assembly analysis:
Exploded view

Inspection:
Cross-section / clipping

Visual analysis:
Realistic · ghosted · blueprint · wireframe

Pipeline:
CAD-derived geometry → optimized real-time assets

Optimization:
6.5M+ triangles → ~95,617 triangles
[verify final metric]

Evaluation:
SUS · NASA-TLX Raw · Think-Aloud

Role:
Unity/C# implementation · technical UI · asset optimization · evaluation

Built as Multimedia Engineering thesis
\end{mdcode}




\medskip\hrule\medskip




\section{11. Silent LinkedIn version}




LinkedIn often autoplays without sound. The silent version must work with captions only.




\subsection{Silent version rules}




\begin{itemize}
\item Use large captions.
\item Keep each caption under 8 words when possible.
\item Use strong visual changes.
\item No tiny UI text.
\item Add final call-to-action.
\item Keep under 60–75 seconds.
\end{itemize}




\subsection{Silent caption sequence}




\begin{mdcode}
Bloque de codigo: text
TwinSight X500

Unity WebGL technical visualization

Problem:
2D assembly diagrams are limited

Goal:
Understand spatial relationships

Interactive drone inspection

Select components

View technical information

Exploded view

Cross-section tools

Visual inspection modes

CAD-to-realtime optimization

Built with:
Unity · C# · WebGL · Blender

Evaluated with:
SUS · NASA-TLX · Think-Aloud

Thesis complete, pending defense

Portfolio / GitHub / LinkedIn
\end{mdcode}




\medskip\hrule\medskip




\section{12. LinkedIn 30–45 second cutdown}




\subsection{Use case}




Use this as a first public LinkedIn post if the full demo is not ready.




\subsection{Structure}




\begin{center}
\scriptsize
\begin{longtable}{>{\raggedright\arraybackslash}p{0.305\linewidth} >{\raggedright\arraybackslash}p{0.305\linewidth} >{\raggedright\arraybackslash}p{0.305\linewidth}}
\toprule
Time & Visual & Caption \\
\midrule
\endfirsthead
\toprule
Time & Visual & Caption \\
\midrule
\endhead
0–3 s & Hero & TwinSight X500 \\
3–8 s & Orbit & Unity WebGL drone inspection \\
8–14 s & Selection & Component selection \\
14–20 s & Exploded & Exploded view \\
20–27 s & Cross-section & Cross-section tools \\
27–34 s & Visual modes & Multiple inspection modes \\
34–42 s & Closing & Unity · C\# · WebGL · Blender \\
\bottomrule
\end{longtable}
\end{center}




\subsection{Caption text for post}




\begin{mdcode}
Bloque de codigo: text
TwinSight X500 is my Unity WebGL technical visualization prototype for drone assembly inspection.

The project explores how interactive real-time 3D can support spatial understanding of technical assembly processes.

Built with Unity, C#, WebGL and Blender.
\end{mdcode}




\medskip\hrule\medskip




\section{13. Technical-reviewer cut}




\subsection{Use case}




Use when applying to technical visualization, simulation, digital twin or Unity technical artist roles.




\subsection{Emphasis}




This version should include more technical footage:




\begin{itemize}
\item Unity editor view;
\item scene hierarchy;
\item scripts briefly;
\item asset before/after;
\item triangle count;
\item WebGL build;
\item interaction architecture diagram;
\item evaluation summary.
\end{itemize}




\subsection{Structure}




\begin{mdcode}
Bloque de codigo: text
1. App overview
2. Interaction systems
3. Asset pipeline
4. WebGL optimization
5. Evaluation
6. Limitations
7. Role/contribution
\end{mdcode}




\subsection{Voiceover focus}




Use more technical phrases:




\begin{mdcode}
Bloque de codigo: text
runtime interaction systems
CAD-derived geometry
WebGL constraints
asset optimization
state management
technical UI
evaluation methodology
\end{mdcode}




Do not overdo code screenshots unless they are readable.




\medskip\hrule\medskip




\section{14. Capture shot checklist}




\subsection{Scene preparation}




\begin{itemize}
\item [ ] Use final or clean build.
\item [ ] Hide debug UI unless needed.
\item [ ] Set good camera starting position.
\item [ ] Use readable UI scaling.
\item [ ] Disable unnecessary console logs.
\item [ ] Use stable frame rate.
\item [ ] Prepare saved camera positions if possible.
\item [ ] Ensure model materials look acceptable.
\item [ ] Prepare visual mode transitions.
\item [ ] Prepare exploded view state.
\item [ ] Prepare cross-section state.
\item [ ] Prepare selected component examples.
\end{itemize}




\subsection{Recording checklist}




\begin{itemize}
\item [ ] Record full drone orbit.
\item [ ] Record component selection.
\item [ ] Record info panel.
\item [ ] Record exploded view activation.
\item [ ] Record exploded view rotation.
\item [ ] Record cross-section/clipping.
\item [ ] Record each visual mode.
\item [ ] Record CAD/optimized comparison if available.
\item [ ] Record metrics/evaluation slide.
\item [ ] Record closing screen.
\end{itemize}




\subsection{Post-capture checklist}




\begin{itemize}
\item [ ] Cut dead time.
\item [ ] Remove errors.
\item [ ] Add captions.
\item [ ] Add title screen.
\item [ ] Add end card.
\item [ ] Normalize audio.
\item [ ] Export 1080p.
\item [ ] Create silent version.
\item [ ] Create thumbnail.
\end{itemize}




\medskip\hrule\medskip




\section{15. Suggested screen recording order}




Record in this order to simplify editing:




\begin{enumerate}
\item full drone idle;
\item camera orbit;
\item component selection;
\item info panel;
\item isolate/highlight behavior;
\item exploded view;
\item exploded view + camera orbit;
\item cross-section/clipping;
\item visual modes;
\item CAD/optimized before-after;
\item metrics/evaluation screen;
\item end card.
\end{enumerate}




\medskip\hrule\medskip




\section{16. Thumbnail options}




\subsection{Thumbnail 1 — best general}




Visual:




\begin{itemize}
\item drone full view;
\item clean dark background;
\item wireframe overlay;
\item short text.
\end{itemize}




Text:




\begin{mdcode}
Bloque de codigo: text
TwinSight X500
Unity WebGL Technical Visualization
\end{mdcode}




\subsection{Thumbnail 2 — technical}




Text:




\begin{mdcode}
Bloque de codigo: text
CAD-to-Realtime Drone Visualization
Unity · WebGL · Blender
\end{mdcode}




\subsection{Thumbnail 3 — recruiter}




Text:




\begin{mdcode}
Bloque de codigo: text
Interactive 3D Drone Assembly Viewer
\end{mdcode}




Avoid thumbnail text longer than 6–8 words.




\medskip\hrule\medskip




\section{17. End card}




\subsection{End card text}




\begin{mdcode}
Bloque de codigo: text
TwinSight X500

Unity · C# · WebGL · Blender
Technical Visualization · CAD-to-Realtime

Portfolio: [URL]
GitHub: [URL]
LinkedIn: [URL]
\end{mdcode}




\subsection{Shorter end card}




\begin{mdcode}
Bloque de codigo: text
TwinSight X500
Unity WebGL Technical Visualization

Portfolio · GitHub · LinkedIn
\end{mdcode}




\medskip\hrule\medskip




\section{18. Music and voice}




\subsection{Voiceover}




Recommended if English pronunciation is acceptable.




If voiceover is weak or noisy, use silent captioned version.




\subsection{Music}




Optional.




If used:




\begin{itemize}
\item low volume;
\item no dramatic cinematic music;
\item no lyrics;
\item no distracting beats;
\item royalty-free;
\item lower than voice.
\end{itemize}




\subsection{Audio priority}




\begin{mdcode}
Bloque de codigo: text
Clear voice > music
\end{mdcode}




If voice is not clear, remove it.




\medskip\hrule\medskip




\section{19. Export formats}




\subsection{Portfolio video}




\begin{mdcode}
Bloque de codigo: text
MP4
1920×1080
30 fps
H.264
90 seconds
\end{mdcode}




\subsection{LinkedIn video}




\begin{mdcode}
Bloque de codigo: text
MP4
1920×1080 or 1080×1080
30 fps
captioned
30–75 seconds
\end{mdcode}




\subsection{GitHub video/GIF}




Use:




\begin{itemize}
\item YouTube/Vimeo/Loom link;
\item or small optimized GIF under \texttt{docs/images/demo-preview.gif}.
\end{itemize}




Avoid large video files inside GitHub repo.




\subsection{File names}




\begin{mdcode}
Bloque de codigo: text
twinsight-x500-demo-90s.mp4
twinsight-x500-recruiter-demo-60s.mp4
twinsight-x500-technical-breakdown-120s.mp4
twinsight-x500-linkedin-cutdown.mp4
\end{mdcode}




\medskip\hrule\medskip




\section{20. Demo page embed copy}




Use this under the embedded video:




\begin{mdcode}
Bloque de codigo: text
TwinSight X500 is a Unity WebGL technical visualization prototype for inspecting drone assembly. The demo shows component selection, exploded view, cross-section tools, visual modes and CAD-to-realtime optimization.
\end{mdcode}




\medskip\hrule\medskip




\section{21. Common mistakes to avoid}




Do not:




\begin{itemize}
\item start with a long thesis explanation;
\item spend 15 seconds on logos;
\item show UI text too small to read;
\item show a laggy or broken WebGL build;
\item overfocus on the drone model beauty;
\item underexplain the technical work;
\item mention every technology in the voiceover;
\item show code that is unreadable;
\item include unverified metrics;
\item use “AI-built” language;
\item show manufacturer assets without disclaimer if public usage is uncertain;
\item exceed 2 minutes for the main demo.
\end{itemize}




\medskip\hrule\medskip




\section{22. Quality gate}




The video is ready when a viewer can answer:




\begin{enumerate}
\item What is TwinSight?
\item What problem does it solve?
\item What did Alexander build?
\item What technologies were used?
\item What makes it relevant to real-time 3D / Unity roles?
\item What features are visible?
\item What evidence exists beyond visuals?
\item Where can I see more?
\end{enumerate}




If the answer to any of these is unclear, revise.




\medskip\hrule\medskip




\section{23. Editing checklist}




\subsection{Before exporting}




\begin{itemize}
\item [ ] Title is clear.
\item [ ] First 5 seconds show the project.
\item [ ] Captions are readable.
\item [ ] Video shows actual interaction.
\item [ ] Feature sequence is logical.
\item [ ] No dead time.
\item [ ] No broken UI.
\item [ ] No unverified metrics.
\item [ ] AI disclosure is not overdone.
\item [ ] End card includes links.
\item [ ] File is under reasonable size.
\item [ ] Audio levels are clean.
\item [ ] Silent version still works.
\end{itemize}




\medskip\hrule\medskip




\section{24. Final recommended production sequence}




\subsection{Step 1}




Record raw footage.




\subsection{Step 2}




Build the 90-second portfolio demo.




\subsection{Step 3}




Cut the 60-second recruiter version.




\subsection{Step 4}




Add captions.




\subsection{Step 5}




Export silent LinkedIn version.




\subsection{Step 6}




Create thumbnail.




\subsection{Step 7}




Embed in portfolio.




\subsection{Step 8}




Add to LinkedIn Featured.




\subsection{Step 9}




Link in GitHub README.




\subsection{Step 10}




Use in outreach.




\medskip\hrule\medskip




\section{25. Next file dependency}




The next file should be:




\begin{mdcode}
Bloque de codigo: text
22_outreach_templates_english_spanish.md
\end{mdcode}




It should include:




\begin{itemize}
\item recruiter connection messages;
\item hiring manager messages;
\item technical director messages;
\item speculative application;
\item follow-up messages;
\item referral requests;
\item salary expectation responses;
\item remote-from-Colombia explanation;
\item contractor setup explanation;
\item AI-assisted development explanation;
\item low formal experience explanation;
\item English and Spanish versions.
\end{itemize}




No Deep Research or agent required.



